\id{ГРНТИ 31.21.19}{}

\begin{header}
\swa{А.Ж.~Мендибаева, С.Д.~Фазылов, С.К.~Кабиева, Р.К.~Жаслан, С.Б.~Жаутикова, А.П.~Богоявленский, О.А.~Нуркенов}{РАЗРАБОТКА ТЕХНОЛОГИИ ПОЛУЧЕНИЯ ВОДОРАСТВОРИМОЙ ФОРМЫ 2-((2-ИЗОНИКОТИНОИЛГИДРАЗОНО)МЕТИЛ)БЕНЗОЙНОЙ КИСЛОТЫ В β-ОЛИГОСАХАРИДНОЙ КАПСУЛЕ}

\tsp{1,2}А.Ж.~Мендибаева,
\tsp{1}С.Д.~Фазылов\envelope,
\tsp{2}С.К.~Кабиева,
\tsp{2}Р.К.~Жаслан,
\tsp{3}С.Б.~Жаутикова,
\tsp{4}А.П.~Богоявленский,
\tsp{1}О.А.~Нуркенов
\end{header}

\begin{affil}
\tsp{1}Институт органического синтеза и углехимии РК, Караганда, Казахстан,

\tsp{2}Карагандинский индустриальный университет, Темиртау, Казахстан,

\tsp{3}Карагандинский медицинский университет, Қараганда, Казахстан,

\tsp{4} Институт микробиологии и вирусологии, Алматы, Казахстан

\corrauthor{Корреспондент-автор: faz8990@mail.ru}
\end{affil}

В статье приведены результаты научных исследований по синтезу и
разработке технологической схемы получения субстанции нового средства -
водорастворимой формы 2-((2-изоникотиноилгид\-разоно)метил)бензойной
кислоты в β-олигосахаридной капсуле, обладающий противовирусной
активностью. Субстанция гинказона получена на основе химической
модификации известного противотуберкулезного препарата изониазида (ГИНК)
в условиях реакции с 2-формилбензойной кислотой. Структура субстрата и
его клатратного комплекса устанавливалась современными
физико-химическими методами ИК- (KBr), ЯМР \tsp{1}Н,
\tsp{13}C-спектроскопии (ДМСО-d\tsb{6} и
CDCl\tsb{3}). Водорастворимый комплекс включения гинказона
получена путем инкапсулирования 2-((2-изоникотиноилгид\-разоно)- метил)
бензойной кислоты с использованием крахмального β-олигосахарида из
водно-спиртовой среды. Проведено молекулярное моделирование механизма
образования клатратных комплексов включения субстрата. Противовирусная
активность исследуемого вещества с целью расчёта его
химико-терапевтического индекса оценивалась в диапазоне концентраций
0,0016-0,2\%, что эквивалентно дозировкам 0,003--0,4 мг на один куриный
эмбрион (в пересчёте 0,06-8 мг/кг массы). Вирусоингибирующая активность
нового соединения в два раза выше, чем у тамифлю и ремантадина. Новое
соединение и его водорастворимый комплекс характеризуются высоким
химико-терапевтическим индексом, что создает перспективу для дальнейшего
использования в фармакологических исследованиях.

{\bfseries Ключевые слова}: гидразид изоникотиновой кислоты,
β-олигосахарид, комплексооброзование, субстанция, антивирусные средства,
технологическая схема.

\begin{header}
2-((2-ИЗОНИКОТИНОИЛГИДРАЗОНО)МЕТИЛ)БЕНЗОЙ ҚЫШҚЫЛЫНЫҢ β-ОЛИГОСАХАРИДПЕН ҚАПТАЛҒАН СУДА ЕРИТІН ТҮРІН АЛУ ТЕХНОЛОГИЯСЫН ӘЗІРЛЕУ

\tsp{1,2}Ә.Ж.~Меңдібаева,
\tsp{1}С.Д.~Фазылов\envelope,
\tsp{2}С.Қ.~Қабиева,
\tsp{2}Р.Қ.~Жаслан,
\tsp{3}С.Б.~Жаутикова,
\tsp{4}А.П.~Богоявленский,
\tsp{1}О.А.~Нүркенов
\end{header}

\begin{affil}
\tsp{1}ҚР органикалық синтез және көмір химиясы институты, Қарағанды, Қазақстан,

\tsp{2}Қарағанды индустриалды университеті, Теміртау, Қазақстан,

\tsp{3}Қарағанды медициналық университеті, Қарағанды, Қазақстан,

\tsp{4}Микробиология және вирусология институты, Алматы, Қазақстан,

e-mail: faz8990@mail.ru
\end{affil}

Бұл мақалада вирусқа қарсы белсенділік көрсететін жаңа субстанция -
2-((2-изоникотиноилгид\-разоно)метил)бензой қышқылының β-олигосахаридпен
қапталған суда еритін түрін алу синтезі және технологиялық схеманы
әзірлеу бойынша ғылыми зерттеулердің нәтижелері ұсынылған. Гинказон заты
туберкулезге қарсы белгілі изониазид (ИНҚГ) препаратын 2-формилбензой
қышқылымен химиялық модификациялау арқылы алынды. Субстраттың құрылымы
және оның клатрат кешені ИҚ- (KBr), \tsp{1}H және
\tsp{13}C ЯМР спектроскопиясы (DMSO-d\tsb{6} және
CDCl\tsb{3}) сияқты заманауи физика-химиялық әдістерді қолдану
арқылы анықталды. Гинказонның суда еритін кешені сулы-спиртті ортада
крахмал β-олигосахаридін пайдаланып 2-((2-изоникотиноилгидразоно)метил)
бензой қышқылын қаптау арқылы алынды. Субстраттың клатрат қосылу
кешендерінің түзілу механизмін молекулалық модельдеу жүргізілді. Қосылыс
үлгісінің вирусқа қарсы белсенділігі тауық эмбрионына 0,003-0,4 мг
(0,06-8 мг/кг) дозада сәйкес келетін 0,0016\%-дан 0,2\%-ға дейінгі
концентрациядағы химиялық-терапиялық индексті анықтау үшін зерттелді.
Жаңа қосылыстың вирусқа қарсы белсенділігі Тамифлю мен римантадинге
қарағанда екі есе жоғары. Жаңа қосылыс және оның суда еритін кешені
жоғары химиялық-терапиялық индексімен сипатталады, бұл оны әрі қарай
фармакологиялық зерттеулер үшін маңызды болып табылады.

{\bfseries Түйін сөздер}: изоникотин қышқылының гидразиді, β-олигоқант,
комплекс түзiлу, гинказон, антивирустық қосылыстар, технологиялық схема.

\begin{header}
DEVELOPMENT OF TECHNOLOGY FOR THE PRODUCTION OF A WATER-SOLUBLE FORM OF 2-((2-ISONICOTINOYLHYDRAZONE)METHYL)BENZOIC ACID IN A β-OLIGOSACCHARIDE CAPSULE

\tsp{1,2}A.Zh.~Mendibayeva,
\tsp{1}S.D.~Fazylov\envelope,
\tsp{2}S.K.~Kabiyeva,
\tsp{2}R.K.~Zhaslan,
\tsp{3}S.B.~Zhautikova,
\tsp{4}A.P.~Bogoyavlensky,
\tsp{1}O.A.~Nurkenov
\end{header}

\begin{affil}
\tsp{1}Institute of organic synthesis and coal chemistry of the Republic of Kazakhstan, Karaganda, Kazakhstan,

\tsp{2}Karaganda industrial university, Temirtau, Kazakhstan,

\tsp{3}Karaganda medical university, Karaganda, Kazakhstan,

\tsp{4}Institute of microbiology and virology, Almaty, Kazakhstan,

e-mail: faz8990@mail.ru
\end{affil}

The article presents the results of scientific research on the synthesis
and development of a technological scheme for the production of a
substance of the new compound a water-soluble form of
2-((2-isonicotinoyl\-hydrazone) methyl)benzoic acid in a β-oligosaccharide
capsule, which has antiviral activity. The gincazon substance was
obtained on the basis of a chemical modification of the well-known
anti-tuberculosis drug isoniazid (IAH) in reaction with 2-formylbenzoic
acid. The structure of the substrate and its clathrate complex was
determined by modern physicochemical methods of IR- (KBr),
\tsp{1}H and \tsp{13}C NMR spectrosco\-py
(DMSO-d\tsb{6} and CDCl\tsb{3}). The water-soluble
gincazon inclusion complex was obtained by encapsulation of
2-((2-isonicotinoylhydrazone) methyl)benzoic acid using starch
β-oligosaccharide from an aqueous alcohol solution. Molecular modeling
of the mechanism of formation of clathrate complexes of substrate
inclusion has been carried out. The viral inhibitory activity of the
compound sample for determining the chemical therapeutic index was
studied at concentrations from 0.0016\% to 0.2\%, which corresponded to
doses of 0.003-0.4 mg per chicken embryo (0.06-8 mg/kg). The
virus-inhibiting activity of the new compound is twice as high as that
of tamiflu and remantadine. The new compound and its water-soluble
complex are characterized by a high chemical-therapeutic index, which
creates prospects for further use in pharmacological research.

{\bfseries Key words}: isonicotinic acid hydrazide, β-oligosaccharide,
antiviral agents, complex formation, ginca\-zon, technological scheme.

\begin{multicols}{2}
{\bfseries Введение}. Согласно данным ВОЗ, ежегодная смертность от
различных форм респираторных заболеваний достигает порядка одного
миллиона случаев. Многократное применение известных противовирусных
препаратов сопровождается формированием резистентности вирусных штаммов
к используемым лекарственным средствам, что существенно снижает
эффективность терапии. Указанные обстоятельства обусловливают
необходимость поиска новых подходов к решению данной проблемы, в том
числе посредством разработки и внедрения новых фармакологически активных
соединений {[}1-3{]}. Значительная часть лекарственных средств,
применяемых в клинической практике для лечения заболеваний органов
дыхания, содержит в своей структуре гидразидный фрагмент. Соединения
данного класса являются достаточно хорошо изученными и находят широкое
применение в различных областях науки, техники и медицины {[}4{]}. В
научной литературе имеется большое количество научных работ, посвящённых
изучению методов получения гидразидных производных, исследованию их
строения и физико-химических свойств. Указанные соединения сохраняют
актуальность как объекты дальнейших исследований и модификации. В связи
с этим использование гидразида изоникотиновой кислоты, как исходного
субстрата в поиске новых антивирусных средств вызывает большой интерес у
химиков-синтетиков и фармакологов всего мира.

В химической практике применяются различные методы перевода
нерастворимых органических субстратов в водорастворимое состояние с
использованием вспомогательных компонентов, например, перевод в солевую
форму (гидрохлориды) {[}5{]}. В настоящее время различные
фармацевтические компании начали широко практиковать метод
капсулирования с использованием натуральных олигосахаридов, например,
β-циклодекстрин (β-ЦД), получаемых из картофельного крахмала. Получение
клатратных комплексов действующих субстратов с циклическими декстринами
будет способствовать повышению их водорастворимости, биодоступности и
химической стабильности {[}6{]}. Кроме того, данный подход может
обеспечивать пролонгирование периода полувыведения активного компонента
и, как следствие, снижение терапевтической дозы препарата. Включение
молекул известных лекарственных средств в комплексы с циклодекстринами
или их производными может приводить к модификации их физико-химических и
фармакокинетических характеристик, что потенциально усиливает
терапевтическую эффективность по сравнению с исходными соединениями.

{\bfseries Материалы и методы.}~Объектами исследования служили
2-((2-изоникотиноилгидразоно) метил) бензойной кислоты, β-циклодекстрин
и их комплекс включения (5). В этом исследований нами приведены новые
научные результаты по получению противовирусного средства (5), а также
описана его технологическая схема получения. Исходное новое соединение
(3) получен методом химической трансформации строения гидразида
изоникотиновой кислоты (ГИНК) (1) с использованием расчетного количества
2-формилбензойной кислоты в спиртовой среде.
\end{multicols}

\fig[0.6\textwidth]{c2/image20}

\begin{multicols}{2}
Инфракрасные спектры соединения (3) регистрировали на спектрометре
ФСМ-1201 с использованием таблеток KBr. Спектры ЯМР \tsp{1}H
и \tsp{13}C получали на приборе JNM-ECA Jeol 400 при рабочих
частотах 399,78 и 100,53 МГц соответственно, в качестве растворителей
применяли ДМСО-d\tsb{6} и CDCl₃. Значения химических сдвигов
(δ, м.д.) в спектрах \tsp{1}H и \tsp{13}C
определяли относительно сигналов остаточных протонов либо атомов
углерода дейтерированных растворителей. Контроль химической чистоты
синтезированных соединений и хода реакций осуществляли методом
тонкослойной хроматографии на пластинах Silufol UV-254 с визуализацией
зон парами йода. Температуры плавления измеряли на приборе SMP10.
Исследование противовирусной активности нового соединения выполнено в
НПЦ «Микробиологии и вирусологии» (г. Алматы) под руководством А. П.
Богоявленским.

Продолжая исследования в этом направлений, с целью перевода в
водорастворимую форму синтезированного нами нового гидразона
изоникотиновой кислоты (3), получен его клатратный комплекс включения
«гость-хозяйн» (2-((2-изоникотиноилгидразоно)метил) бензойная
кислота:β-ЦД) с β-олигосахаридом (5) в соотношений 1:2 в водно-спиртовой
среде.
\end{multicols}

\fig{c2/image10}[Рис. 1 - Схема образования капсулированной формы соединения (5)]

\begin{multicols}{2}
Далее была разработана лабораторная технологическая схема получения
новой противовирусной субстанций «Гинказон».

\emph{Стадия 1.} Взвешивание гидразида изоникотиновой кислоты и
2-формибензойной кислоты. Взвешивают 2,74 г (гидразида изоникотиновой
кислоты), 3 г 2-формибензойной кислоты передают на стадию 2.

Подготовка растворителя. В мерный цилиндр приливают 50 мл этанола и
передают на стадию 2.

\emph{Стадия 2.} Синтез (2-((2-изоникотиноилгидра\-зоно)метил) бензойной
кислоты.

\emph{Стадия 3.} Реакционную массу переносили в чашку Петри и
выдерживали на воздухе до испарения основной части растворителя.
Оставшийся растворитель удаляли путём упаривания под пониженным
давлением с использованием роторного испарителя.

Осаждение реакционной смеси. Реакционную смесь охлаждают, при этом
выпадает осадок белого цвета.

\emph{Стадия 4.} Фильтрация осадка. Образовавшийся при охлаждении осадок
отделяют фильтрованием и дважды промывают гексаном, сушат и получают
субстанцию (2-((2-изоникотиноилгидразоно)метил) бензойной кислоты.
Полученный технический продукт передают на стадию 5.

\emph{Стадия 5,6.} Обработка фильтрата. Очистку продукта производят в
изопропиловом спирте путем перекристаллизации. Остаток маточника
выпаривают на роторном испарителе.

\emph{Стадия 7.} Для выделения субстанции растворитель отгоняют в
вакууме, остаток хроматографируют на колонке с силикагелем (элюент -
хлороформ, смесь хлороформа с этанолом (100:1→10:1).

\emph{Стадия 8.} Получение водорастворимой формы
(2-((2-изоникотиноилгидразоно)- метил) бензойной кислоты с
β-циклодекстрином - (5). К раствору 0,067 г (0,00025 М)
(2-((2-изоникотиноилгидразоно)метил) бензойной кислоты в 10 мл этанола
при перемешивании добавляют 0,56 г (0,0005 М) β-циклодекстрина в 10 мл
воды. Реакционную массу выдерживают при перемешивании в интервале
температур 50-65°С в течение 3 часов. Выход (5) 63 \%, т. пл.
275-283\tsp{о}С. Белый порошок.

\emph{Стадия 9.} Фильтрация осадка. Осадок, образовавшийся в процессе
охлаждения, подвергают фильтрации, сушат и получают водорастворимую
форму (5). Полученный комплекс передают на стадию 10.

\emph{Стадия 10.} Реакционную массу переносят в чашку Петри и
выдерживают на воздухе до испарения основной части растворителя.
Оставшийся растворитель удаляют путём упаривания под вакуумом с
использованием роторного испарителя.

Полученное соединение (5) отфильтровывают, промывают и сушат. На первом
этапе субстанцию (5) высушивают на фильтровальной бумаге, на втором
этапе - в вакуумном сушильном шкафу до постоянной массы (температура
50-60°С).

\emph{Стадия 11.} Готовый продукт. Субстанция (5) должен иметь вид
белого порошка, без запаха. Субстанцию взвешивают на весах и
упаковывают.
\end{multicols}

\fig[0.85\textwidth]{c2/image21}[Рис.2 - Схематическое описание стадии получения субстанции (5)]

\begin{multicols}{2}
{\bfseries Обсуждение и результаты.}

Методика получения 2-((2-изоникотиноилгид\-разоно)метил) бензойной кислоты
(3). К смеси 2,74 г (0,02 М) гидразида изоникотиновой кислоты в 20 мл
этанола добавляют при перемешивании 3 г (0,02 М) 2-карбоксибензальдегида
в 30 мл этанола. Реакционную смесь перемешивают при температуре
60\tsp{о}С в течение 2 часов, а затем охлаждают до комнатной
температуры. Продукт отфильтровают, промывают и сушат.
Перекристаллизация из изопропанола дала (3) в виде белого порошка. Выход
92 \%, т. пл.205-206\tsp{о}С.

ИК спектр (KBr), ν, см\tsp{-1}: 3425 (NH), 1655 (COамид),
1608 (C=N), 1547 (аром.). Спектр ЯМР \tsp{1}Н
(ДМСО-d\tsb{6}), δ, м.д., (\emph{J}, Гц): 7.49-7.51 (1Н, м,
Н\tsp{15}), 7.60-7.62 (1Н, м, Н\tsp{13}), 7.80
(2Н, с, Н\tsp{3,5}), 7.86-7.89 (1Н, м,
Н\tsp{14}), 8.02-8.05 (1Н, м, Н\tsp{16}), 8.73
(2Н, с, Н\tsp{2,6}), 9.18 (1Н, д, Н\tsp{11},
\tsp{3}\emph{J} = 5.2), 12.25 (1Н, с, Н\tsp{9}),
12.08 (1Н, с, Н\tsp{20}). Спектр ЯМР \tsp{13}С
(ДМСО-d\tsb{6}), δ\tsb{С}, м.д.: 127.27
(С\tsp{16}), 130.43 (С\tsp{17}), 130.88
(С\tsp{15}), 131.33 (С\tsp{13}), 134.88
(С\tsp{12}), 122.14 (С\tsp{3,5}), 148.36
(С\tsp{4}), 150.83 (С\tsp{2,6}), 148.36
(С\tsp{11}), 162.34 (С\tsp{7}), 168.55
(С\tsp{18}). Спектр ЯМР COSY:
Н\tsp{15}→Н\tsp{13},
Н\tsp{15}→Н\tsp{14},
Н\tsp{13}→Н\tsp{16},
Н\tsp{3,5}→Н\tsp{2,6}. Спектр ЯМР HMQC:
Н\tsp{3,5}→С\tsp{3,5},
Н\tsp{15}→С\tsp{15},
Н\tsp{13}→С\tsp{13},
Н\tsp{14}→С\tsp{14},
Н\tsp{16}→С\tsp{16},
Н\tsp{2,6}→С\tsp{2,6},
Н\tsp{11}→С\tsp{11}. Спектр ЯМР HMBC:
Н\tsp{3,5}→С\tsp{2,6},
Н\tsp{3,5}→С\tsp{7};
Н\tsp{2,6}→С\tsp{3,5},
Н\tsp{2,6}→С\tsp{4};
Н\tsp{11}→С\tsp{12},
Н\tsp{11}→С\tsp{16};
Н\tsp{19}→С\tsp{7},
Н\tsp{19}→С\tsp{11}. Масс-спектр, m/z,
(\emph{I}\tsb{отн},\%): 270.142 (100) {[}M+H{]}+.
\end{multicols}

\fig[0.6\textwidth]{c2/image11}
\fig[0.6\textwidth]{c2/image12}[Рис.3 - \tsp{1}Н, \tsp{13}С ЯМР
спектры 2-((2-изоникотиноилгидразоно)метил)бензойной кислоты (3)]

\begin{multicols}{2}
\emph{Определение вирусоингибирующей активности.} Для оценки
вирусингибирующей активности исходного субстрата (3) были использованы
штаммы вируса гриппа с различными антигенными характеристиками:
A/Алматы/8/98 (H3N2) и A/Владивосток/2/09 (H1N1). Исследование
вирусингибирующего эффекта образцов соединения (3) проводилось для
определения химического терапевтического индекса (ХТИ) при концентрациях
от 0,0016\% до 0,2\%, что соответствовало дозам 0,003-0,4 мг на куриный
эмбрион (0,06--8 мг/кг) (табл.1). Из анализа полученных данных следует,
что соединение (3) характеризуются высоким ХТ-индексом, что создает
перспективу его для дальнейшего исследования. Установлено, что
соединение (3) не проявляет токсические свойства. Вирусоингибирующая
активность нового соединения (3) - в два раза выше, чем у тамифлю и
ремантадина.

В таблице 1 приведена вирусоингибирующая активность нового гидразона (3)
в сравнении с эталонными препаратами тамифлю и ремантадином.
\end{multicols}

\tcap{Таблица 1 - Вирусингибирующая активность соединения (3) по отношению к вирусам гриппа}
\begin{longtblr}[
  label = none,
  entry = none,
]{
  cells = {c},
  cells = {font = \small},
  row{1} = {font=\bfseries\small},
  cell{1}{1} = {r=2}{},
  cell{1}{2} = {c=2}{},
  hlines,
  vlines,
}
Соединение и внутренний шифр & Химический терапевтический индекс препаратов & \\
 & А/Алматы/8/98 (H3N2) & А/Владивосток/2/09 (H1N1)\\
(3) & 65,0 & 70,0\\
Тамифлю & 29,9 & 30\\
Ремантадин & 10.3 & 11
\end{longtblr}

\begin{multicols}{2}
\emph{Молекулярное моделирование механизма образования клатратного
комплекса включения.} Исследование комплекса включения для системы
«циклодекстрин-гидразон» в стехиометрии 1:1 позволило оценить афинность
связывания двух реагирующих компонентов и сравнить эффективность
образования комплекса. В качестве гостевой молекулы использовали
гидразон (3), а в роли молекулярного рецептора - β-ЦД (рис.3)
{[}7-9{]}.
\end{multicols}

\fig{c2/image13}[Рис.3 - Схематическое представление комплексообразования
гидразона изоникотиновой кислоты (3) с β-циклодекстрином]

\begin{multicols}{2}
Для подготовки молекулярных структур {[}10-11{]} к процедуре докинга их
геометрия была предварительно оптимизирована методом DFT (RB3LYP/6-31G,
модель растворителя CPCM), что позволило определить наиболее стабильные
конформации и проанализировать их структурные характеристики. На рисунке
4 представлены оптимизированные геометрии энергетически предпочтительных
конформаций гидразона и β-ЦД при образовании комплекса включения.
\end{multicols}

\fig{c2/image14}[Рис.4 - Оптимизированные геометрии объектов исследования]

\begin{multicols}{2}
{\bfseries Выводы.} В результате выполненных исследований на основе
химической модификации известного противотуберкулезного препарата
изониазида получена его олигосахаридная композиция. Водорастворимый
комплекс включения гинказона получена путем инкапсулирования
2-((2-изоникотиноилгидразоно)метил) бензойной кислоты с использованием
β-олигосахарида. Проведено молекулярное моделирование механизма
образования клатратных комплексов включения субстрата. Противовирусная
активность исследуемого вещества с целью расчёта его ХТИ оценивалась в
диапазоне концентраций 0,0016--0,2\%, что эквивалентно дозировкам
0,003--0,4 мг на один куриный эмбрион (в пересчёте 0,06--8 мг/кг массы).

Вирусоингибирующая активность нового соединения в два раза выше, чем у
тамифлю и ремантадина. Новое соединение и его водорастворимый комплекс
характеризуются высоким химико-терапевтическим индексом, что создает
перспективу его для дальнейших фармакологических исследований.
Полученные результаты являются научной основой для разработки
технологической схемы получения нового антивирусного средства с высокой
активностью.
\end{multicols}

\begin{center}
{\bfseries Литература}
\end{center}

\begin{refs}
1. Wolker N., Howe C., Glover M., McRobbie H., Barnes J., et al.
Cytisine versus Nikotine for Smoking Cessation // The New England
Journal of Medicine. - 2014. - Vol.371(25). - Р.
2353-2362.~\href{https://doi.org/10.1056/nejmoa1407764}{DOI
10.1056/nejmoa1407764}.

2. Thomas D., Farrel M., McRobbie H., Tutka P., et al. The
effectiveness, safety and cost‐effectiveness of cytisine versus
varenicline for smoking cessation in an Australian population: A study
protocol for a randomized controlled non‐inferiority trial // Addiction.
-2018. -Vol.114(5). - P.923-933.
\href{https://doi.org/10.1111/add.14541}{DOI 10.1111/add.14541}.

3. Fedorova V.A., Kadyrova R.A., Slita A.V., Muryleva A.A., Petrova
P.R., Kovalskaya A.V., Lobov A.N., Zileeva Z.R., Tsypyshev D.O.,
Borisevich S.S., Tsypysheva I.P., Vakhitova J.V., Zarubaev V.V.
Antiviral activity of amides and carboxamides of quinolizidine alkaloid
(-)-cytisine against human influenza virus A(H1N1) and parainfluenza
virus type 3 // Natural Product Research. -2021. -Vol.35(22).
-Р.4256-4264.~\href{https://doi.org/10.1080/14786419.2019.1696791}{DOI
10.1080/14786419.2019.1696791}.

4. Prochaska J.J., Das S., Benowitz N.L. Cytisine, the world's oldest
smoking cessation aid // BMJ. ­- 2013. - Vol.347: f5198. DOI
10.1136/bmj.f5198.

5. Vakhitova Yu.V., Farabontova E.I., Zainullina L.F., Vakhitov V.A.,
Tsypysheva I., Yunusov M.S. Search of (-)-cytisine derivatives as
potential inhibitors of NF-kB and STAT1 // Russian Journal of Biorganic
chemistry. -2015. -Vol.41(3). -P.297-304.
DOI~\href{http://dx.doi.org/10.1134/S1068162015030103}{10.1134/S1068162015030103}.

6. Tsypysheva I.P., Koval'skaya A., Petrova P., Lobov A., Borisevich
S.S., Tsypyshev D., Fedorova V.A., Gorbunova E.A., Galochkina A.V.,
Zarubaev V.V. Diels-Alder Adducts of Nsubstituted Derivatives of
(-)-Cytisine as Influenza A/H1N1 Virus Inhibitors; Stereo
differentiation of antiviral Properties and Preliminary Assessment of
Action Mechanism // Tetrahedron.- 2019. -Vol.75(21). - P.2933-2943.
DOI
\href{https://doi.org/10.1016/j.tet.2019.04.021}{10.1016/j.tet.2019.04.021}.

7. Nurkenov O.A., Fazylov S.D., Satpaeva Zh.B., Seilkhanov T.M.,
Turdybekov D.M., Mendibayeva A.Zh., Akhmetova S.B., Shulgau Z.T.,
Alkhimova L.E., Kulakov I.V. Synthesis, structure and biological
activity of hydrazones derived from 2- and 4-hydroxybenzoic acid
hydrazides //Chemical Data Collections. -2023. -Vol.48: 01089.
\href{https://doi.org/10.1016/j.cdc.2023.101089}{DOI10.1016/j.cdc.2023.101089}.

8. Muldakhmetov Z., Fazylov S., Gazaliev A., Nurkenov O., Seilkhanov O.
The synthesis of new inclusion compounds complexes
cytosine-β-cyclodextrin // News of the national Academy of Sciences of
the Republic of Kazakhstan Series Chemistry and technology. -2022. -Vol.
2(451). - Р.112-120.
\href{https://doi.org/10.32014/2022.2518-1491.107}{DOI
10.32014/2022.2518-1491.107}.

9. Larsen K.L. Large cyclodextrins // Journal of Inclusion Phenomena and
Macrocyclic Chemistry. -2002. -Vol.43(1). -P.1-13.
\href{https://doi.org/10.1023/A:1020494503684}{DOI
10.1023/A:1020494503684}.

10. Perkin E. Informatics (Computer software) // ChemOffice
Professional. - 2019. (Version 19.0). URL:
\url{https://www.perkinelmer.com/category/chemdraw}. -Дата обращения:
16.01.2026.

11. Revvity Signals Software, Inc.~ChemDraw Professional~(Version 22.2).
-2023.
URL:~\url{https://shop.revvitysignals.com/product/chemdraw-professional}.
-Дата обращения: 16.01.2026.
\end{refs}

\begin{info}
\hspace{1em}\emph{{\bfseries Сведения об авторах}}

Мендибаева А.Ж. - младший научный сотрудник, Институт органического
синтеза и углехимии РК, Караганда, Казахстан, e-mail:
anenyawa@mail.ru;

Фазылов С.Д. - академик НАН РК, доктор химических наук, Институт
органического синтеза и углехимии РК, Караганда, Казахстан, e-mail:
faz8990@mail.ru;

Кабиева С.К. - кандидат химических наук, Карагандинский индустриальный
университет, Темиртау, Казахстан, e-mail: kabieva.s@mail.ru;

Жаслан Р.К. - PhD, доцент, Карагандинский индустриальный университет,
Темиртау, Казахстан, e-mail: rymgul.zhaslan@mail.ru;

Жаутикова С.Б. - академик НАЕН РК, доктор медицинских наук,
Карагандинский медицинский университет, Караганда, Казахстан, e-mail:
zhautikova@qmu.kz;

Богоявленский А.П. - академик РАЕ, доктор биологических наук,
профессор, РГП «Институт микробиологии и вирусологии», Алматы,
Казахстан, e-mail: anpav\_63@mail.ru;

Нуркенов О.А. - доктор химических наук, заведующий лабораторией СБАВ,
Институт органического синтеза и углехимии РК, Караганда, Казахстан,
e-mail: nurkenov\_oral@mail.ru.

\hspace{1em}\emph{{\bfseries Information about the authors}}

Mendibayeva A.Zh. - Junior Researcher, Institute of Organic Synthesis
and Coal Chemistry of the Republic of Kazakhstan, Karaganda,
Kazakhstan; e-mail: anenyawa@mail.ru;

Fazylov S.D. - Academician of the National Academy of Sciences of the
Republic of Kazakhstan, Doctor of Chemical Sciences, Institute of
Organic Synthesis and Coal Chemistry of the Republic of Kazakhstan,
Karaganda, Kazakhstan, e-mail: faz8990@mail.ru;

Kabieva S.K. - Candidate of Chemical Sciences, Karaganda Industrial
University, Temirtau, Kazakhstan, e-mail: kabieva.s@mail.ru;

Zhaslan R.K. - PhD, Associate Professor, Karaganda Industrial
University, Temirtau, Kazakhstan, e-mail: rymgul.zhaslan@mail.ru;

Zhautikova S.B. - Academician of the Kazakhstan National Academy of
Natural Sciences, Doctor of Medical Sciences, Karaganda Medical
University, Karaganda, Kazakhstan, e-mail: zhautikova@qmu.kz;

Bogoyavlensky A.P. - Academician of the Russian Academy of Natural
Sciences, Doctor of Biological Sciences, Professor, The RSE "Institute
of Microbiology and Virology", Almaty, Kazakhstan, e-mail:
anpav\_63@mail.ru;

Nurkenov O.A. - Doctor of Chemical Sciences, Institute of Organic
Synthesis and Coal Chemistry of the Republic of Kazakhstan, Karaganda,
Kazakhstan, e-mail: nurkenov\_oral@mail.ru.
\end{info}
